\documentclass[12]{article}
\usepackage{graphicx}
\usepackage[T1]{fontenc}
\usepackage{tgbonum}
\usepackage{amsmath}
\usepackage{amssymb}
\usepackage{booktabs}
\usepackage{float}
\usepackage{graphicx} % Required for inserting images
\usepackage{ulem}
\usepackage{hyperref}
\usepackage{fontawesome5}
\usepackage{dirtytalk}
\usepackage{xcolor}
\usepackage[affil-it]{authblk}

\usepackage{setspace}
\doublespacing

\usepackage{geometry}
\geometry{
 a4paper,
 total={170mm,257mm},
 left=20mm,
 top=20mm,
 }






\title{EK369E: Financial Markets\\ Bond Prices and Yields\\ Problem Set}
\author{Nord University Business School}
\date{}

\begin{document}

\maketitle

\color{black}
\section*{Problem 1}
Consider an 6\% coupon bond selling for NOK 93.3 with three years to maturity making annual coupon payments. This bond is redeemable at maturity at par value.
\begin{enumerate}
    \item[a)] Calculate the bond’s yield to maturity (YTM). 
    \item[b)] If interest rates over the next three years will be, with certainty, $r_1 = 6\%$, $r_2 = 8\%$ and $r_3 = 10\%$, calculate the realised compound average return on this bond.
\end{enumerate}





\color{black}
\section*{Problem 2}
A 20 Year maturity bond with par value of NOK 100 makes annual coupon payments at 8\%. Find the yield to maturity of the bond if the bond's price is
\begin{enumerate}
    \item[a)] NOK 95
    \item[b)] NOK 100
    \item[c)] NOK 105
\end{enumerate}



\color{black}
\section*{Problem 3}
Fill in the table below for the following zero-coupon bonds with par values of NOK 100.

\begin{table}[H]
    \centering
    \resizebox{0.5\textwidth}{!}{%
    \begin{tabular}{ccc}
    \toprule
        Price & Years to Maturity & Yield To Maturity (YTM)\\\midrule
        NOK 40 & 20 & a) \\
        NOK 50 & 20 & b)\\
        NOK 50 & 10 & c)\\
        d) & 10 & 10\%\\
        e) & 10 & 8\%\\
        NOK 40 & f) & 8\%\\\bottomrule
    \end{tabular}
    }
    \label{tab:my_label}
\end{table}





\color{black}
\section*{Problem 4}
Bond Alfa has a current yield of 9\% and a yield to maturity of 10\%. 
\begin{enumerate}
    \item[a)] Is Bond Alfa selling above or below par value?
    \item[b)] Is the coupon rate on Bond Alfa more or less than 9\%?
\end{enumerate}




\color{black}
\section*{Problem 5}
A newly issued bond pays its coupons once annually. Its coupon rate is 5\%, its maturity is 20 years and its yield to maturity is 8\%.
\begin{enumerate}
    \item[a)] Find the holding period return for a 1-year investment period if the bond is selling at a yield to maturity of 7\% by the end of the year.
    \item[b)] Find the realised compound yield for a two-year holding period, assuming that you sell the bond after two years, the bond's yield is 7\% at the end of the second year and coupons can be reinvested for one year at 3\% interest rate.
\end{enumerate}




\color{black}
\section*{Problem 6}
The table below summarises some information from the fixed-income market:
% Please add the following required packages to your document preamble:
% \usepackage{booktabs}
% \usepackage{graphicx}
\begin{table}[H]
\centering
\resizebox{0.4\textwidth}{!}{%
\begin{tabular}{@{}llll@{}}
\toprule
Bond & \begin{tabular}[c]{@{}l@{}}Maturity\\ (years)\end{tabular} & \begin{tabular}[c]{@{}l@{}}Coupon rate\\ (\%)\end{tabular} & \begin{tabular}[c]{@{}l@{}}Yield\\ (\%)\end{tabular} \\ \midrule
A    & 5.5                                                        & 0                                                          & 2.5                                                  \\
B    & 2.5                                                        & 4.0                                                        & 2.0                                                  \\ \bottomrule
\end{tabular}%
}
\label{tab:my-table}
\end{table}

\begin{enumerate}
    \item[a)] Calculate the invoice price and the flat price for the two bonds.
    \item[b)] How can a zero-coupon bond be traded above its face value?
\end{enumerate}




\color{black}
\section*{Problem 7}
Treasury bonds paying an 8\% coupon rate with semiannual payments currently sell at par value. what coupon rate would they have to pay in order to sell at par if they paid their coupons annually. 




\color{black}
\section*{Problem 8}
Assume that you have a 1-year investment horizon and are trying to choose among three bonds with the same degree of default risk and the same time to maturity of 10 years. Bond 1 is a zero coupon bond, Bond 2 has a coupon rate of 8\%, and Bond 3 has a coupon rate of 10\%
\begin{enumerate}
    \item[a)] If all bonds are currently priced to yield 8\% at maturity, what are the prices of Bond 1, Bond 2 and Bond 3
    \item[b)] If at the beginning of next year, all bonds are priced to yield 8\% , what will the prices of these bonds be?
    \item[c)] Based on your answer in a) and b) compute the holding period return over the 1 year horizon for each bond. 
    \item[d)] recalculate b) and c) if at the beginning of next year, all bonds are priced to yield 7\%
\end{enumerate}
\end{document}


